\begin{workedexample}[Worked Example: A slot as a leakage antenna]
A shield's effectiveness is set less by wall thickness than by its
apertures. A slot radiates efficiently once its length approaches a half
wavelength; leakage becomes severe near $f \approx c/(2L)$. A $30\ \text{cm}$ seam
leaks strongly around
\[ f = \frac{3\times10^8}{2\times0.30} = 500\ \text{MHz}. \]
The rule: keep the longest slot dimension $\ll \lambda/2$ --- many small holes
beat one long slot of equal area.
\end{workedexample}

\begin{keytakeaways}
\begin{itemize}[leftmargin=1.4em,itemsep=2pt]
\item Shielding effectiveness $=$ reflection $+$ absorption loss; apertures dominate real enclosures.
\item A slot leaks strongly as its length nears $\lambda/2$ ($f \approx c/2L$); slot length, not area, governs leakage.
\item Cable shields couple via transfer impedance; pigtails ruin high-frequency shielding.
\end{itemize}
\end{keytakeaways}

\begin{exercises}
\item At what frequency does a $15\ \text{cm}$ slot become a strong radiator?
\item For fixed total open area, are many small holes or one long slot better?
\item What ruins a cable shield's HF performance at the connector?
\end{exercises}

\furtherreading{Ott~\cite{ott} (ch.~6); Paul~\cite{paul}.}
